\section{Introduction}
\label{sec:intro}

Transistor network synthesis is a foundational step in the design and
customization of CMOS standard cells~\cite{asap7,bonncell,nvcell}.
Given a Boolean function $f$, the goal is to build a minimum-transistor
switching network—an undirected graph whose edges represent transistors
and whose source-to-sink connectivity implements $f$.
The synthesized topology directly determines transistor count,
diffusion-sharing opportunities, and downstream layout
quality~\cite{switchcraft,bonncell};
even a single extra transistor can degrade area, power, and timing
after placement and routing, so automatic cell-layout tools depend on a
compact transistor network at their input.

\paragraph{Limitations of existing methods.}
Transistor-network synthesis ranges from heuristic switch-graph
construction~\cite{kagaris2007,switchcraft,motox} to exact SAT-based
solvers~\cite{muSTNet,miniTNtk} that can provably find
minimum-transistor solutions.
All of them, however, solve each function from scratch,
with no reuse across instances.
Runtime grows exponentially with the number of variables, and the methods
offer no generalization to unseen functions.
None of these methods \emph{learns} a generalizable synthesis policy,
so every new function incurs the full cost of a fresh optimization.

\paragraph{Our approach: learned sequential graph generation.}
We treat transistor network synthesis as a sequential
graph-generation Markov decision process (MDP) and train a policy
to solve it.
At each step the policy places one transistor, guided by a graph neural
network that observes the current partial circuit together with a
compact representation of the target function.
A single policy trained over a distribution of Boolean functions
generalizes at inference time to new functions without any per-instance
search.
Learning has recently reshaped many stages of the design
flow~\cite{mlforeda}, including reinforcement-learning-driven standard-cell
layout~\cite{nvcell}, yet transistor-network synthesis itself has remained
the domain of per-instance exact solvers.

\paragraph{Key technical ideas.}
\textbf{(1) Theoretical guarantees by construction.}

We prove a \emph{Monotonicity} theorem showing that adding transistors
never removes existing conducting paths, and derive a
\emph{Safety Inheritance} corollary that reduces safety verification
at each step from checking the whole network to a single BFS on the
new edge.
Safety and correctness are therefore maintained as invariants throughout
the entire trajectory, not checked post-hoc.

\textbf{(2) A novel GNN architecture for transistor synthesis.}
We design \method, whose core is the \sgpn (\SGPN) — a four-head
factored policy built on a Message Passing Neural Network (MPNN)~\cite{gilmer2017mpnn} backbone.
Each head specializes in one sub-decision of the action
(source node, drain node, control variable, polarity),
with safety masks applied independently at each head.
A learnable virtual-node embedding allows the policy to propose
fresh topology without pre-committing to a node count.
All heads share a single MPNN forward pass over a
\emph{unified graph} that merges the partial synthesized network
with a fixed SOP compensate network.

\textbf{(3) Three-phase training toward minimum transistor count and physical quality.}
Phase~1 NLL pretraining on known-optimal solutions initializes the policy.
Phase~2 PPO finetuning with a GRPO group-relative baseline, an adaptive
self-tightening reward, and an entropy bonus drives convergence toward
minimum transistor count without mode collapse.
Phase~3 extends the RL loop to joint PDN+PUN synthesis with a physical
placement reward from ASTRAN~\cite{astran}, directly optimizing the
weighted placement objective (cell width, gate mismatch, wire length,
routing density, and diffusion gaps).

\paragraph{Contributions.}
\begin{enumerate}[leftmargin=*, topsep=2pt, itemsep=3pt]

  \item \textbf{An MDP formulation that is correct and safe by construction.}
        We cast transistor network synthesis as a safety-masked sequential
        graph-generation MDP and prove a Monotonicity theorem (adding
        transistors is order-invariant for conducting paths) with a Safety
        Inheritance corollary that reduces safety checking to
        $O(|\Poff|\cdot|V|)$ per candidate, independent of how many transistors
        are already placed.  Together they guarantee that every masked
        trajectory is \emph{correct and safe by construction}, with no post-hoc
        verification (Section~\ref{sec:theory}).

  \item \textbf{The four-head \SGPN{} architecture.}
        \method{} factors each transistor placement into four masked
        sub-decisions over a \emph{unified graph} that fuses the partial circuit
        $G_t$ with a fixed SOP compensate network $C_g$ encoding the target, so
        a shared MPNN sees both current topology and remaining coverage.  A
        virtual new-node embedding lets the policy open internal nodes without a
        node-count decision, per-head safety masks make every sampled action
        safe, and a batched \texttt{PackedGraph} forward pass replaces
        $O(K\cdot T_{\max})$ MPNN calls with one (Section~\ref{sec:method}).

  \item \textbf{Three-phase training and exhaustive evaluation.}
        A single policy is driven from imitation toward minimum count and
        layout quality: NLL pretraining, then GRPO-based transistor-count RL
        with a self-tightening reward, then a physical-aware phase that jointly
        synthesizes the PDN ($\lnot f$) and PUN ($f(\lnot\mathbf{x})$) and
        rewards them by ASTRAN's~\cite{astran} placement cost.  On the
        exhaustive three-variable sweep \method{} reaches the exact minimum
        count on all $508$ targets ($97.6\%$ already after pretraining), extends
        to the four- and six-variable sweeps, and lowers mean placement cost by
        $18\%$/$15\%$ over count-optimal cells---with no per-instance solver
        call at inference (Sections~\ref{sec:method}--\ref{sec:exp}).

\end{enumerate}

The remainder of the paper is organized as follows.
Section~\ref{sec:pre} introduces transistor network synthesis and
placement/routing background.
Section~\ref{sec:method} presents our theoretical foundations and the
\method{} architecture with three-phase training.
Section~\ref{sec:exp} reports experimental results and ablations.
Section~\ref{sec:con} concludes.
